\section{From Legal Labels to Economic Type}
\label{sec:cobidder_type}

The validation section treats cobidders as the loser-side
cartel-adjacency target. That is the right legal object for an
award-layer screen, but it leaves an economic question. Are
cobidders simply frequent losers who happened to appear near CADE
defendants, or do they look like the type of firm a cover-bidding
screen is meant to find? We read this question within the
frequent-loser stratum, comparing cobidders with frequent losers
that do not co-bid with direct defendants. This is the hard
comparison: both groups lose persistently and both cross the same
participation threshold.

\subsection{Firm-Level Signature}
\label{sec:cobidder_firm}

Cobidders differ from other frequent losers in the expected
directions on most firm-level dimensions. They participate more
often, face more unique winners, bid more often in tenders that
include a direct CADE defendant, and concentrate activity in
narrower item-group portfolios. The unique-winner contrast is
especially large: cobidders face $\valBridgeUniqWinCob$ unique
winners on average, compared with $\valBridgeUniqWinFLnc$ for
frequent-loser non-cobidders (Cohen's
$d=\valBridgeUniqWinD$). The proximity measure also moves in the
right direction: $\valBridgeDirectCob$ of cobidders' loss-bids
face a direct CADE defendant, compared with
$\valBridgeDirectFLnc$ for other frequent losers
($d=\valBridgeDirectD$). Their portfolios are more concentrated
as well, with item-group HHI $\valBridgeHHICob$ against
$\valBridgeHHIFLnc$ ($d=\valBridgeHHID$).

One dimension goes the wrong way. Cobidders have a lower share of
own loss-bids in repeat winner-pair clusters than frequent-loser
non-cobidders ($\valBridgeRepeatShareCob$ versus
$\valBridgeRepeatShareFLnc$). We do not use that measure as
support for the cover-bidding interpretation. The broader
firm-level profile is still informative: the cobidder label is
not just a proximity tag attached to otherwise identical active
losers.

\subsection{Bid-Level Signature}
\label{sec:cobidder_bid}

The LANCES export lets us check whether the same distinction
appears inside bids. For firms with at least five usable losing
bids in convite or pregão, we compute the gap between each losing
bid and the winning bid in the same tender-item,
$(\text{bid}-\text{winning bid})/\text{winning bid}$, winsorized
at $[-0.99,10]$ before firm-level aggregation. The usable sample
contains $\valTBBLnCob$ cobidders and $\valTBBLnFLnc$
frequent-loser non-cobidders.

The bid-level pattern is not the old textbook version of cover
bidding, in which cover bids are plainly uncompetitive. Cobidders
bid closer to winners, not farther away: their median bid gap is
$\valTBBLMedGapCob$, compared with $\valTBBLMedGapFLnc$ for
other frequent losers ($d=\valTBBLMedGapD$). At the same time,
their within-firm dispersion in bid gaps is higher
($\valTBBLSDGapCob$ versus $\valTBBLSDGapFLnc$,
$d=\valTBBLSDGapD$). A logit of cobidder status within the
frequent-loser stratum, controlling for
$\log(1+\text{tenders\_count})$, preserves both signs: the
median-gap coefficient is $\valTBBLLogitMedCoef$
($z=\valTBBLLogitMedZ$), and the dispersion coefficient is
$\valTBBLLogitSDCoef$ ($z=\valTBBLLogitSDZ$). The bid-level
distinction is therefore not only another way of measuring high
participation.

\begin{table}[htbp]
\centering
\caption{Cobidder signature within the frequent-loser stratum}
\label{tab:cobidder_signature_v18}
\footnotesize
\begin{adjustbox}{max width=\textwidth,center}
\begin{tabular}{p{4.4cm}ccc}
\toprule
Dimension & Cobidders & FL non-cobidders & Difference \\
\midrule
Total participations & $\valBridgeTendCob$ & $\valBridgeTendFLnc$ & $d=\valBridgeTendD$ \\
Unique winners faced & $\valBridgeUniqWinCob$ & $\valBridgeUniqWinFLnc$ & $d=\valBridgeUniqWinD$ \\
Loss-bids facing direct CADE defendant & $\valBridgeDirectCob$ & $\valBridgeDirectFLnc$ & $d=\valBridgeDirectD$ \\
Portfolio item-group HHI & $\valBridgeHHICob$ & $\valBridgeHHIFLnc$ & $d=\valBridgeHHID$ \\
Repeat winner-pair share & $\valBridgeRepeatShareCob$ & $\valBridgeRepeatShareFLnc$ & $d=\valBridgeRepeatShareD$ \\
\midrule
Median bid gap to winner & $\valTBBLMedGapCob$ & $\valTBBLMedGapFLnc$ & $d=\valTBBLMedGapD$ \\
Within-firm SD of bid gap & $\valTBBLSDGapCob$ & $\valTBBLSDGapFLnc$ & $d=\valTBBLSDGapD$ \\
\bottomrule
\end{tabular}
\end{adjustbox}
\end{table}

\subsection{What the Signature Establishes}
\label{sec:cobidder_interpretation}

The pattern is most consistent with credible cover bidding.
Cover bids must lose, but they need not be absurdly high. In
many cartels they must remain plausible enough to preserve the
appearance of rivalry and to avoid drawing attention
\citep{asker2010leniency,marshall2012economics}. A firm that
cycles through cover roles may therefore bid relatively close to
the winner in some tenders while showing elevated dispersion
across tenders. That is the pattern we observe.

This is not a proof of cartel membership. The LANCES evidence is
used for type validation, not adjudication. It narrows the most
important alternative reading of the validation results: that
cobidders are merely high-volume losers mechanically adjacent to
CADE defendants. They are high-volume losers, but they also have
a firm-level and bid-level profile that fits the economic type
the award-layer screen is meant to rank. With that target in
place, the next question is whether the award-layer screen can
organize the more expensive bid-layer tools used in cartel
forensics.
